\section{Prediction of Stationary Processes (Chapter 5)}
\label{sec:tmts-ch5}
Unless otherwise stated, $X$ is a real, centered stationary process with
covariance $\gamma$. Write
\[
\mathcal H_n=\operatorname{span}\{X_1,\ldots,X_n\},\qquad
\mathcal M_n=\overline{\operatorname{span}}\{X_t:t\le n\}.
\]
$P_{\mathcal H}$ always denotes orthogonal projection in $L^2$.
For a nonzero known mean $\mu$, apply the results to $X_t-\mu$ and add
$\mu$ to the predicted value.

\subsection{Prediction Equations and Optimality (Section 5.1)}
\begin{definition}[Best linear predictor]
The best linear predictor of $X_{n+h}$ from $X_1,\ldots,X_n$ is
\[
\widehat X_{n+h\mid n}=P_{\mathcal H_n}X_{n+h}.
\]
It uniquely minimizes $\mathbb E(X_{n+h}-Y)^2$ over $Y\in\mathcal H_n$
and is characterized by
\[
\mathbb E[(X_{n+h}-\widehat X_{n+h\mid n})Y]=0,
\qquad Y\in\mathcal H_n.
\]
\end{definition}
\begin{theorem}[Normal equations]
Let $\Gamma_n=[\gamma(i-j)]_{i,j=1}^n$ and
$\boldsymbol\gamma_n^{(h)}=(\gamma(h),\ldots,\gamma(h+n-1))^\top$.
Then
\[
\widehat X_{n+h\mid n}=\sum_{j=1}^n\phi_{nj}^{(h)}X_{n+1-j},\qquad
\Gamma_n\boldsymbol\phi_n^{(h)}=\boldsymbol\gamma_n^{(h)}.
\]
A solution always exists. The predictor is unique even when its coefficient
vector is not. If $\Gamma_n$ is nonsingular, then
\[
\boldsymbol\phi_n^{(h)}=\Gamma_n^{-1}\boldsymbol\gamma_n^{(h)},\qquad
v_n(h)=\gamma(0)-(\boldsymbol\gamma_n^{(h)})^\top
\Gamma_n^{-1}\boldsymbol\gamma_n^{(h)}.
\]
Here $v_n(h)=\mathbb E(X_{n+h}-\widehat X_{n+h\mid n})^2$.
\end{theorem}
\begin{proposition}[Nonsingularity; Proposition 5.1.1]
If $\gamma(0)>0$ and $\gamma(h)\to0$ as $|h|\to\infty$, then
$\Gamma_n$ is nonsingular for every $n$. This is sufficient, not necessary.
\end{proposition}
For every $Y\in\mathcal H_n$, the orthogonal decomposition gives
\[
\mathbb E(X_{n+h}-Y)^2=v_n(h)+
\mathbb E(\widehat X_{n+h\mid n}-Y)^2.
\]
Thus the normal equations express an optimization property, rather than
merely a procedure for finding coefficients.

\subsection{Recursive Prediction and Innovations (Section 5.2)}
\begin{theorem}[Durbin--Levinson recursion; Proposition 5.2.1]
Suppose all $\Gamma_n$ are positive definite. For one-step prediction write
$\phi_{nj}=\phi_{nj}^{(1)}$ and $v_n=v_n(1)$, with $v_0=\gamma(0)$.
Then $\phi_{11}=\gamma(1)/\gamma(0)$ and, for $n\ge2$,
\[
\kappa_n:=\phi_{nn}=
\frac{\gamma(n)-\sum_{j=1}^{n-1}\phi_{n-1,j}\gamma(n-j)}{v_{n-1}},
\]
\[
\phi_{nj}=\phi_{n-1,j}-\kappa_n\phi_{n-1,n-j},
\qquad 1\le j<n,
\]
\[
v_n=v_{n-1}(1-\kappa_n^2),\qquad n\ge1,
\quad \kappa_1=\phi_{11}.
\]
The reflection coefficient $\kappa_n$ is the partial autocorrelation
$\alpha(n)$. Positive definiteness implies $|\kappa_n|<1$.
\end{theorem}
\begin{corollary}[Prediction variance and determinant]
\[
v_n=\gamma(0)\prod_{j=1}^n(1-\alpha(j)^2),\qquad
\det\Gamma_n=\prod_{j=0}^{n-1}v_j.
\]
The decreasing sequence $v_n$ is the error remaining after increasingly
long observation histories.
\end{corollary}
\begin{definition}[Finite-history innovations]
Set $\widehat X_{1\mid0}=0$ and
\[
e_t=X_t-\widehat X_{t\mid t-1},\qquad
\mathbb E e_t^2=v_{t-1}.
\]
The $e_t$ are pairwise orthogonal and
$\mathcal H_n=\operatorname{span}\{e_1,\ldots,e_n\}$.
They need not be independent or have equal variances.
\end{definition}
\begin{theorem}[Innovations representation and recursion]
Assume $v_j>0$. The one-step predictor can be written
\[
\widehat X_{n+1\mid n}=\sum_{j=1}^n\vartheta_{nj}e_{n+1-j}.
\]
With $v_0=\gamma(0)$, the coefficients satisfy, for $0\le k<n$,
\[
\vartheta_{n,n-k}=\frac{1}{v_k}
\left(\gamma(n-k)-\sum_{j=0}^{k-1}
\vartheta_{k,k-j}\vartheta_{n,n-j}v_j\right),
\]
\[
v_n=\gamma(0)-\sum_{j=0}^{n-1}\vartheta_{n,n-j}^2v_j.
\]
An empty sum is zero. Unlike Durbin--Levinson, which updates coefficients
of observations, this recursion updates coefficients of orthogonal innovations.
\end{theorem}
The innovations recursion also applies to a centered nonstationary sequence.
Writing $K(s,t)=\mathbb E(X_sX_t)$ and assuming every finite covariance
matrix is nonsingular, its general form is
\[
\vartheta_{n,n-k}=v_k^{-1}\left(K(n+1,k+1)-
\sum_{j=0}^{k-1}\vartheta_{k,k-j}\vartheta_{n,n-j}v_j\right),
\]
\[
v_0=K(1,1),\qquad
v_n=K(n+1,n+1)-\sum_{j=0}^{n-1}\vartheta_{n,n-j}^2v_j.
\]
Stationarity reduces $K(n+1,k+1)$ to $\gamma(n-k)$.
\begin{proposition}[Multiple-step prediction in innovations coordinates]
For $h\ge1$,
\[
\widehat X_{n+h\mid n}=\sum_{j=1}^n
\frac{\mathbb E(X_{n+h}e_j)}{v_{j-1}}e_j,
\qquad
v_n(h)=\gamma(0)-\sum_{j=1}^n
\frac{[\mathbb E(X_{n+h}e_j)]^2}{v_{j-1}}.
\]
Writing $\mathbf X_n=(X_1,\ldots,X_n)^\top$ and
$X_t=\sum_{j=1}^t L_{tj}e_j$, with $L_{tt}=1$, gives the factorization
\[
\Gamma_n=L\operatorname{diag}(v_0,\ldots,v_{n-1})L^\top.
\]
\end{proposition}

\subsection{Finite-History ARMA Prediction (Section 5.3)}
Let $X$ be causal ARMA$(p,q)$, and put
$\widehat X_{s\mid n}=X_s$ for $1\le s\le n$.
\begin{proposition}[ARMA prediction equations]
For $n+h>\max(p,q)$, projection of the ARMA equation gives
\[
\widehat X_{n+h\mid n}
=\sum_{i=1}^p\phi_i\widehat X_{n+h-i\mid n}
+\sum_{j=h}^q\theta_jP_{\mathcal H_n}Z_{n+h-j}.
\]
In particular, if $h>q$, the noise term vanishes and the forecasts satisfy
the homogeneous AR difference equation. Finite-history projections of
$Z_s$ need not equal $Z_s$, even when the model is invertible.
\end{proposition}
\begin{proposition}[ARMA innovations form; Section 5.3]
Put $m=\max(p,q)$ and introduce the triangularly transformed observations
\[
Y_t=\begin{cases}X_t,&1\le t\le m,\\
\phi(B)X_t,&t>m.
\end{cases}
\]
For every $n$, $\operatorname{span}\{Y_1,\ldots,Y_n\}=\mathcal H_n$,
and the finite-history innovations of $Y$ are precisely $e_t$.
Apply the general innovations recursion to the covariance of $Y$, writing
its coefficients as $\eta_{nj}$. For $n\ge m$,
\[
\widehat X_{n+1\mid n}=\sum_{i=1}^p\phi_iX_{n+1-i}
+\sum_{j=1}^q\eta_{nj}e_{n+1-j}.
\]
For $n+h>m$, the multiple-step version is
\[
\widehat X_{n+h\mid n}=\sum_{i=1}^p\phi_i\widehat X_{n+h-i\mid n}
+\sum_{j=h}^{\min(q,n+h-1)}\eta_{n+h-1,j}e_{n+h-j}.
\]
Here $\widehat X_{s\mid n}=X_s$ for observed indices. The coefficients
$\eta_{nj}$ refer to $Y$, and generally differ from the coefficients
$\vartheta_{nj}$ obtained by applying the recursion directly to $X$.
\end{proposition}
For arbitrary horizons, the normal equations give the exact finite-history
error covariance. With
$\mathbf c_h=(\gamma(n+h-1),\ldots,\gamma(h))^\top$, it is
\[
\operatorname{Cov}(X_{n+h}-\widehat X_{n+h\mid n},
X_{n+k}-\widehat X_{n+k\mid n})
=\gamma(h-k)-\mathbf c_h^\top\Gamma_n^{-1}\mathbf c_k.
\]
This formula applies to all stationary processes with nonsingular
$\Gamma_n$, including ARMA models.

\subsection{Gaussian Prediction and Prediction Bounds (Section 5.4)}
\begin{theorem}[Conditional Gaussian law]
For a centered stationary Gaussian process with nonsingular $\Gamma_n$,
\[
\mathbb E[X_{n+h}\mid X_1,\ldots,X_n]=\widehat X_{n+h\mid n}.
\]
The prediction error is independent of $(X_1,\ldots,X_n)$ and normally
distributed with variance $v_n(h)$. If $v_n(h)>0$, then conditionally on
the observed history,
\[
X_{n+h}\sim N(\widehat X_{n+h\mid n},v_n(h)).
\]
Hence, for $0<\beta<1$, an exact conditional $(1-\beta)$ prediction interval is
\[
\widehat X_{n+h\mid n}\ \pm\ z_{1-\beta/2}\sqrt{v_n(h)}.
\]
These statements assume the model parameters are known.
\end{theorem}
For a general non-Gaussian process, the best linear predictor need not equal
the conditional expectation, which minimizes mean squared error over all
square-integrable functions of the observations.

\subsection{Prediction from the Infinite Past (Section 5.5)}
\begin{theorem}[Causal invertible ARMA prediction; Theorem 5.5.1]
Suppose
\[
X_t=\sum_{j=0}^{\infty}\psi_jZ_{t-j},\qquad
Z_t=\sum_{j=0}^{\infty}\pi_jX_{t-j},\qquad\psi_0=\pi_0=1,
\]
with both coefficient sequences absolutely summable. Then
$\mathcal M_n=\overline{\operatorname{span}}\{Z_t:t\le n\}$, and
\[
\widetilde X_{n+h\mid n}:=P_{\mathcal M_n}X_{n+h}
=\sum_{j=h}^{\infty}\psi_jZ_{n+h-j}.
\]
The error and its variance are
\[
E_{n,h}=\sum_{j=0}^{h-1}\psi_jZ_{n+h-j},\qquad
v_\infty(h)=\sigma^2\sum_{j=0}^{h-1}\psi_j^2.
\]
Writing $\widetilde X_{s\mid n}=X_s$ for $s\le n$ also gives
\[
\widetilde X_{n+h\mid n}=-\sum_{j=1}^{\infty}\pi_j
\widetilde X_{n+h-j\mid n}.
\]
\end{theorem}
\begin{proposition}[Joint forecast errors]
For $h,k\ge1$,
\[
\operatorname{Cov}(E_{n,h},E_{n,k})
=\sigma^2\sum_{\ell=1}^{\min(h,k)}\psi_{h-\ell}\psi_{k-\ell}.
\]
Thus forecast errors at different horizons generally are correlated.
Moreover $v_\infty(h)\uparrow\gamma(0)$ as $h\to\infty$.
\end{proposition}
Increasing finite histories ending at a fixed time yield projections that
converge in $L^2$ to $P_{\mathcal M_n}X_{n+h}$.
By stationarity, $v_n(h)\downarrow v_\infty(h)$ as the number of observed
past values tends to infinity, for each fixed $h$.

\subsection{Prediction in the Frequency Domain (Section 5.6)}
\begin{theorem}[Projection under the spectral isometry]
Let $I(g)=\int g\,dZ$ be the spectral isometry and set
\[
\mathcal K_n=\overline{\operatorname{span}}^{L^2(F)}
\{e^{it\lambda}:t\le n\}.
\]
Then
\[
P_{\mathcal M_n}X_{n+h}=I(P_{\mathcal K_n}e^{i(n+h)\lambda}),
\]
\[
\mathbb E|X_{n+h}-P_{\mathcal M_n}X_{n+h}|^2
=\inf_{g\in\mathcal K_n}\int|e^{i(n+h)\lambda}-g(\lambda)|^2\,dF(\lambda).
\]
Time-domain prediction is therefore equivalent to weighted $L^2$
approximation of exponentials by past-frequency functions.
\end{theorem}
For a causal invertible ARMA process, another expression for the predictor is
\[
P_{\mathcal M_n}X_{n+h}=\sum_{j=0}^{\infty}a_jX_{n-j},\qquad
\sum_{j=0}^{\infty}a_jz^j=
\frac{\phi(z)}{\theta(z)}\sum_{j=0}^{\infty}\psi_{j+h}z^j.
\]

\subsection{The Wold Decomposition (Section 5.7)}
\begin{definition}[Determinism and pure nondeterminism]
Let $\mathcal M=\overline{\operatorname{span}}\{X_t:t\in\mathbb Z\}$ and
$\mathcal M_{-\infty}=\bigcap_{n\in\mathbb Z}\mathcal M_n$.
The process is deterministic if
\[
\sigma_\varepsilon^2:=\mathbb E(X_t-P_{\mathcal M_{t-1}}X_t)^2=0,
\]
equivalently $X_t\in\mathcal M_{-\infty}$ for every $t$.
It is purely nondeterministic if $\mathcal M_{-\infty}=\{0\}$.
These are statements about linear predictability, not about sample paths
being nonrandom.
\end{definition}
\begin{theorem}[Wold; Theorem 5.7.1]
If $\sigma_\varepsilon^2>0$, there is a unique decomposition
\[
X_t=U_t+V_t,\qquad U_t=\sum_{j=0}^{\infty}\psi_j\varepsilon_{t-j},
\]
where the series converges in mean square and
\[
\psi_0=1,\quad\sum_{j=0}^{\infty}\psi_j^2<\infty,\qquad
\varepsilon_t=X_t-P_{\mathcal M_{t-1}}X_t
\sim\mathrm{WN}(0,\sigma_\varepsilon^2).
\]
The coefficients and deterministic component are
\[
\psi_j=\sigma_\varepsilon^{-2}\mathbb E(X_t\varepsilon_{t-j}),
\qquad V_t=P_{\mathcal M_{-\infty}}X_t.
\]
$\{V_t\}$ is stationary and deterministic, and
$\mathbb E(V_s\varepsilon_t)=0$ for all $s,t$.
Furthermore,
\[
\mathcal M_t=
\overline{\operatorname{span}}\{\varepsilon_s:s\le t\}
\mathbin{\oplus}\mathcal M_{-\infty},\qquad
\overline{\operatorname{span}}\{V_s:s\le t\}=\mathcal M_{-\infty}.
\]
\end{theorem}
Wold coefficients need only be square summable. The representation is
therefore more general than the absolutely summable $\mathrm{MA}(\infty)$
definition in Chapter 3.
\begin{corollary}[Prediction from the Wold representation]
\[
P_{\mathcal M_t}X_{t+h}
=V_{t+h}+\sum_{j=h}^{\infty}\psi_j\varepsilon_{t+h-j},
\qquad v_\infty(h)=\sigma_\varepsilon^2\sum_{j=0}^{h-1}\psi_j^2.
\]
The deterministic part contributes no prediction error.
As $h\to\infty$, the error variance tends to $\operatorname{Var}(U_t)$.
\end{corollary}
\begin{theorem}[Spectral decomposition; Theorem 5.7.2]
Under $\sigma_\varepsilon^2>0$, the spectral measure decomposes as
$F_X=F_U+F_V$, where $F_U$ is absolutely continuous and $F_V$ is singular.
The density of $U$ is
\[
f_U(\lambda)=\frac{\sigma_\varepsilon^2}{2\pi}
\left|\sum_{j=0}^{\infty}\psi_je^{-ij\lambda}\right|^2.
\]
The Fourier series is interpreted as an $L^2$ boundary function.
The assumption $\sigma_\varepsilon^2>0$ is essential: a deterministic
process may have an absolutely continuous spectral measure.
\end{theorem}

\subsection{Kolmogorov's Formula (Section 5.8)}
\begin{theorem}[Kolmogorov; Theorem 5.8.1]
Let $f$ be the density of the absolutely continuous part of the spectral
measure of a real centered stationary process. Its infinite-past one-step
prediction error variance is
\[
\boxed{\displaystyle
\sigma_\varepsilon^2=2\pi\exp\left\{
\frac1{2\pi}\int_{-\pi}^{\pi}\log f(\lambda)\,d\lambda\right\}.}
\]
The convention is $\exp(-\infty)=0$. Consequently,
\[
\sigma_\varepsilon^2>0
\quad\Longleftrightarrow\quad
\int_{-\pi}^{\pi}\log f(\lambda)\,d\lambda>-\infty.
\]
The singular part of the spectral measure does not enter this formula.
\end{theorem}
\begin{corollary}[Pure nondeterminism in spectral terms]
A nonzero stationary process is purely nondeterministic precisely when
its spectral measure is absolutely continuous and its density satisfies
$\int_{-\pi}^{\pi}\log f(\lambda)\,d\lambda>-\infty$.
Positivity of $f$ almost everywhere alone is insufficient; logarithmic
integrability is also required.
\end{corollary}
